\documentclass[quiet]{ctexart}

\ExplSyntaxOn \makeatletter

\keys_define:nn { hduthesis / exam }
  {
    watermark .tl_set:N = \l__hduthesis_exam_watermark_tl,
  }
\NewDocumentCommand \hdutexset {m} { \keys_set:nn { hduthesis / exam } {#1} }
\DeclareHookRule { shipout / background }
  { hduthesis / exam } { before } { pgfrcs }
\AddToHook { shipout / background } [ hduthesis / exam ]
  {
    \tl_if_empty:NF \l__hduthesis_exam_watermark_tl
      {
        \tikz [ remember~picture, overlay ]
          \node [ opacity = .2 ] at (current~page)
            {
              \includegraphics [ width = \linewidth ]
                { \l__hduthesis_exam_watermark_tl }
            };
      }
  }

\RequirePackage
  {
    geometry, mathtools, array, bm, esint, keytheorems, libertine, inconsolata,
    siunitx, fixdif, derivative, physics2, tasks, anyfontsize
  }
\RequirePackage[inline]{enumitem}
\RequirePackage{fontawesome5, fancyqr, circuitikz}
\hook_gput_code:nnn { cmd/faIcon/before } { . }
  {
    \group_begin:
    \fontsize { \fp_eval:n { .9 * \f@size } } { 0 } \selectfont
  }
\hook_gput_code:nnn { cmd/faIcon/after } { . } { \group_end: }
\geometry{margin = 1in}

\RequirePackage[svgnames]{xcolor}
\newkeytheorem{note}
  [ style = remark, numbered = false,
    name = \color{DarkRed}\sffamily \faIcon{info-circle}~Tip. ]
\newkeytheorem{problem}
  [ style = definition,
    name = \color{DarkGreen}\sffamily \faIcon{pen-square}~Problem ]
\newkeytheorem{solution}
  [ style = definition, numbered = false,
    name = \color{MidnightBlue}\sffamily \faIcon{check-square}~Solution ]
\usephysicsmodule{ab, op.legacy}
\usephysicsmodule{bm-um.legacy}
\settasks { label = \sffamily\Alph*., item-indent = 2.5em, label-offset = .5pt }
\newcommand \true {\addtocounter{task}{1}\faIcon{check-circle}}
\def \@result#1 {\underline{~#1~}}
\def \s@result {\textsf{【\qquad 】}}
\DeclareRobustCommand \result {\@ifstar\s@result\@result}
\keys_define:nn { textfig }
  {
    ratio      .tl_set:N    = \l__textfig_ratio_tl,
      ratio    .initial:n   = 7 \c_colon_str 2,
    position   .tl_set:N    = \l__textfig_position_tl,
      position .initial:n   = c,
    unknown    .code:n      = \tl_if_novalue:nF {#1}
      { \tl_set_eq:NN \l__textfig_position_tl \l_keys_key_tl }
  }
\NewDocumentEnvironment{textfig}{ O{} }
  {
    \group_begin:
    \keys_set:nn { textfig } {#1}
    \seq_set_split:Nee \l__whats_textfig_ratio_seq { \c_colon_str }
      { \l__textfig_ratio_tl }
    \noindent
    \minipage [ \l__textfig_position_tl ]
      {
        \fp_eval:n
          {
            .96*\seq_item:Nn \l__whats_textfig_ratio_seq {1}/
              (
                \seq_item:Nn \l__whats_textfig_ratio_seq {1} + 
                \seq_item:Nn \l__whats_textfig_ratio_seq {2}
              )
          }\dimexpr \linewidth
      }
  }
  {
    \endminipage
    \vspace{4.5pt plus 1.5pt minus 1.5pt} \par
    \group_end:
  }
\NewDocumentCommand \textfigsplit {}
  {
    \endminipage
    \hspace{.02\linewidth}
    \minipage [ \l__textfig_position_tl ]
      {
        \fp_eval:n
          {
            .96*\seq_item:Nn \l__whats_textfig_ratio_seq {2}/
              (
                \seq_item:Nn \l__whats_textfig_ratio_seq {1} + 
                \seq_item:Nn \l__whats_textfig_ratio_seq {2}
              )
          }\dimexpr \linewidth
      }
  }

% \hook_gput_code:nnn { cmd/maketitle/before } { . } { \vskip -\baselineskip }
% \hook_gput_code:nnn { cmd/maketitle/after } { . }
%   { \__weyoung_maketitle_aftercode_aux: }
\cs_new_protected:Nn \__weyoung_maketitle_aftercode_aux:
  {
    \vspace*{-2\baselineskip}
    \begin{center}
      \begin{minipage}{.66\textwidth}
        \centering
        \begin{tasks}
          [ label-width = 5\ccwd, item-indent = 2cm,
            after-item-skip = 0pt ](2)
          \task[\textbf{课程信息}：] 理学院/A0715022
          \task[\textbf{授课教师}：] 陈许敏
          \task[\textbf{解析提供}：] \textsf{Skyrmion}
          \task[{\makebox[5em][s]{\textbf{\textsf\TeX~排版}：}}] \textsf{Skyrmion}
        \end{tasks}
      \end{minipage}
      \begin{minipage}{.15\textwidth}
        \centering
        \begin{tikzpicture}
          \node [ inner~sep = 0 ] (wechat) {
            \fancyqr
              [ size = 4em, image~padding = 1,
                image = \raisebox{4em/33}{\textcolor{DarkGreen}{\faIcon{weixin}}},
                r~color = DarkGreen, version = 4
              ]{http://weixin.qq.com/r/sRJiemTEGjllrSAZ90ej}
            };
          \node [ inner~sep = .5ex, align = center, right, font = \scriptsize ]
            at (wechat.east) {未\\央\\学\\社};
        \end{tikzpicture}
      \end{minipage}
      \begin{minipage}{.15\textwidth}
        \centering
        \begin{tikzpicture}
          \node [ inner~sep = 0 ] (qq) {
            \fancyqr
              [ size = 4em, image~x~padding = 2, image~y~padding = 1,
                image = \raisebox{4em/33}{\textcolor{MidnightBlue}{\faIcon{qq}}},
                l~color = MidnightBlue, version = 4
              ]{https://qm.qq.com/q/PH7CsQWvm2}
            };
          \node [ inner~sep = .5ex, align = center, left, font = \scriptsize ]
            at (qq.west) {物\\理\\营};
        \end{tikzpicture}
      \end{minipage}
    \end{center}
  }

\newcommand \iu {{\mathup i\mkern1mu}}
\newcommand \eu {\mathup e}
\AtBeginDocument{%
  \let \Re \relax \let \Im \relax
  \DeclareMathOperator{\Re}{Re}
  \DeclareMathOperator{\Im}{Im}
}

\makeatother \ExplSyntaxOff
\usetikzlibrary{ fadings, patterns, arrows, decorations.markings }
\tikzset
  {
    > = stealth',
    every picture/.append style =
      { line join = round, line cap = round, thick }
  }

\setCJKmainfont{LXGW WenKai}[AutoFakeSlant]
\setCJKsansfont{LXGW Marker Gothic}[AutoFakeBold, AutoFakeSlant, Scale = .95]
\setCJKmonofont{LXGW WenKai Mono}[AutoFakeBold, AutoFakeSlant]
\ctexset{section/format+ = \raggedright}
\usepackage[warnings-off={mathtools-colon,mathtools-overbracket}]{unicode-math}
\setmathfont[AutoFakeBold]{STIX Two Math}

\makeatletter \def \@result #1{\underline{ \quad }} \makeatother
\DeclareMathOperator{\DFT}{DFT}

\title{\sffamily\bfseries 24Fall \textsc{Hdu} 「高等数学 A1」期末试题}
\author{} \date{}

\begin{document}

\maketitle

\section{选择题（每题 3 分，共 24 分）}

\begin{problem}
  极限 $\lim_{x \to \infty} (\sqrt{x^2 + x} - x) =$
  \begin{tasks}(4)
    \task $0$
    \task $\frac12$
    \task $\infty$
    \task 不存在
  \end{tasks}
\end{problem}
\begin{solution}

\end{solution}

\begin{problem}
  曲线 $y = \frac{x^2 + x}{x^2 - 1}$ 的渐近线条数为
  \begin{tasks}(4)
    \task $0$
    \task $1$
    \task $2$
    \task $3$
  \end{tasks}
\end{problem}
\begin{solution}
  
\end{solution}

\begin{problem}
  已知 $f(x)$ 的一个原函数为 $\ln x$，则 $f'(x) =$
  \begin{tasks}(4)
    \task $\frac1x$
    \task $-\frac1{x^2}$
    \task $x\ln x - x + C$
    \task $e^{x^2}$
  \end{tasks}
\end{problem}
\begin{solution}
  
\end{solution}

\begin{problem}
  设 $\lim_{x \to 0} \frac{f(x) - f(0)}{1 - \cos x} = 1$，则 $f(x)$ 在 $x = 0$ 处
  \begin{tasks}(4)
    \task 可导且 $f'(0) \neq 0$
    \task 取极小值
    \task 取极大值
    \task 不可导
  \end{tasks}
\end{problem}
\begin{solution}
  
\end{solution}

\begin{problem}
  设函数 $f(x) =
  \begin{cases}
    xe^{x^2}, & -1 \leq x < 1\\
    -1,       & x \geq 1
  \end{cases}$，
  则 $\int_0^3 f(x - 1) \d x =$
  \begin{tasks}(4)
    \task $-1$
    \task $1$
    \task $-2$
    \task $2$
  \end{tasks}
\end{problem}
\begin{solution}
  
\end{solution}

\begin{problem}
  以下反常积分收敛的是
  \begin{tasks}(4)
    \task $\int_e^{+\infty} \frac{\ln x}{x} \d x$
    \task $\int_e^{+\infty} \frac{1}{x \ln x} \d x$
    \task $\int_e^{+\infty} \frac{1}{x (\ln x)^2} \d x$
    \task $\int_e^{+\infty} \frac{1}{x \sqrt{\ln x}} \d x$
  \end{tasks}
\end{problem}
\begin{solution}
  
\end{solution}

\begin{problem}
  设函数 $f(x)$ 连续，则 $\odv{}x \int_0^x tf(x^2 - t^2) \d t =$
  \begin{tasks}(4)
    \task $xf(x^2)$
    \task $-xf(x^2)$
    \task $2xf(x^2)$
    \task $-2xf(x^2)$
  \end{tasks}
\end{problem}
\begin{solution}
  
\end{solution}

\begin{problem}
  若线性无关函数 $y_1$，$y_2$，$y_3$ 是二阶非齐次线性方程
  $y'' + P(x)y' + Q(x)y = f(x)$ 的三个特解，$C_1$，$C_2$ 且为任意常数，
  则此方程的通解为
  \begin{tasks}(2)
    \task $C_1y_1 + C_2y_2 + y_3$
    \task $C_1y_1 + C_2y_2 - (C_1 + C_2)y_3$
    \task $C_1y_1 + C_2y_2 - (1 - C_1 - C_2)y_3$
    \task $C_1y_1 + C_2y_2 + (1 - C_1 - C_2)y_3$
  \end{tasks}
\end{problem}
\begin{solution}
  
\end{solution}

\section{填空题（每题 3 分，共 12 分）}

\begin{problem}
  设 $y = \ln \cos x$，则 $\d y =$ \result{}.
\end{problem}
\begin{solution}
  
\end{solution}

\begin{problem}
  $\int_{-1}^1 \ln (x + \sqrt{1 + x^2}) \d x =$ \result{}.
\end{problem}
\begin{solution}
  
\end{solution}

\begin{problem}
  曲线 $y = \int_0^x \sqrt{\sin t} \d t$（$0 \leq x \leq \pi$）的弧长为 \result{}.
\end{problem}
\begin{solution}
  
\end{solution}

\begin{problem}
  微分方程 $(y + x^2 e^x) \d x - x \d y = 0$ 的通解为 \result{}.
\end{problem}
\begin{solution}
  
\end{solution}

\section{计算题（每题 6 分，共 24 分）}

\begin{problem}
  隐函数 $y = y(x)$ 由方程 $x^3 - \int_0^y e^{-t^2} \d t + y^3 + 4$ 所确定，
  求 $\odv yx$.
\end{problem}
\begin{solution}
  
\end{solution}

\begin{problem}
  计算 $\int \ln (1 + x^2) \d x$.
\end{problem}
\begin{solution}
  
\end{solution}

\begin{problem}
  计算定积分 $\int_0^{\pi/2} e^{2x} \cos x \d x$.
\end{problem}
\begin{solution}
  
\end{solution}

\begin{problem}
  求极限 $\lim_{x \to 0}
  \frac{\int_0^x [3 sin t + t^2 + \cos(1/t)] \d t}{\int_0^x \ln(1 + t) \d t}$.
\end{problem}
\begin{solution}
  
\end{solution}

\section{解答题（共 40 分）}

\begin{problem}[7 pt]
  已知 $\int_0^{+\infty} \frac{\sin x}{x} \d x = \frac\pi2$，
  求 $\int_0^{+\infty} \frac{\sin^2 x}{x^2} \d x$.
\end{problem}
\begin{solution}
  
\end{solution}

\begin{problem}[7 pt]
  已知函数 $f(x)$ 具有二阶导数，若曲线 $y = f(x)$ 过点 $(0,0)$，
  且与曲线 $y = 2^x$ 在点 $(1,2)$ 处相切，求 $\int_0^1 xf''(x) \d x$.
\end{problem}
\begin{solution}
  
\end{solution}

\begin{problem}[7 pt]
  求函数 $y = \frac{x^3}{(x - 1)^2}$ 的极值，凹凸区间和拐点.
\end{problem}
\begin{solution}
  
\end{solution}

\begin{problem}[7 pt]
  求微分方程 $y'' + 3y' + 2y = 3xe^{-x}$ 的通解.
\end{problem}
\begin{solution}
  
\end{solution}

\begin{problem}[8 pt]
  已知抛物线 $y = ax^2$ （$a > 0$） 将两坐标轴和抛物线 $y = 4 - x^2$（$x \geq 0$
  所围成的屏幕区域 $D$ 分成面积相等的左、右两部分区域 $D_1$ 和 $D_2$.
  \begin{enumerate}
    \item 求 $a$ 的值.
    \item 分别求区域 $D_1$ 和 $D_2$ 绕 $x$ 轴旋转一周所得的旋转体的体积.
  \end{enumerate}
\end{problem}
\begin{solution}
  
\end{solution}

\begin{problem}[4 pt]
  设 $f(x)$ 在区间 $[0,1]$ 上连续，在 $(0,1)$ 内可导，且 $f(x) > 0$，
  且 $2\int_{1/2}^1 f(x) \d x = f(0)$. 证明：在 $(0,1)$ 内至少存在一点 $\xi$，
  使 $f'(\xi) = 0$.
\end{problem}
\begin{solution}
  \begin{proof}
    
  \end{proof}
\end{solution}

\end{document}